\documentclass[10pt,a4paper]{amsart}
\usepackage[margin=19mm]{geometry}
\usepackage{amsmath,amssymb,mathtools}
\usepackage[scr=rsfs,cal=euler]{mathalfa}
\usepackage{xfrac}
\usepackage[unicode,hidelinks,bookmarks=false]{hyperref}
\newcommand{\ie}{i.e.\ }
\newcommand{\cf}{cf.\ }
\newcommand{\resp}{respectively\ }
\newcommand{\Subsection}{\subsection}
\providecommand{\nobreakdash}{\nobreak\hbox{-}}
\setlength{\emergencystretch}{2em}
\multlinegap0pt
\usepackage[shortlabels]{enumitem}
\usepackage{amssymb}
\usepackage{mathtools}
\newtheorem{lemma}{Lemma}
\newcommand{\D}{\mathbb D}
\newcommand{\Hh}{\mathbb H}
\newcommand{\R}{\mathbb R}
\title{The Hayman--Wu constant is $\pi^2$}
\author{Paata Ivanisvili}
\date{Source: arXiv:2608.12844v2 (CC BY 4.0)}
\begin{document}
\maketitle

\noindent\emph{Source and licence.} The Hayman--Wu constant is $\pi^2$. Version arXiv:2608.12844v2 (CC BY 4.0), distributed by arXiv under the Creative Commons Attribution 4.0 International licence, http://creativecommons.org/licenses/by/4.0/, as recorded in the arXiv licence metadata for this exact version and shown on the abstract page https://arxiv.org/abs/2608.12844v2. Full licence notice: \texttt{licenses/arxiv-2608.12844-CC-BY-4.0.txt}.\par\smallskip

\noindent\textit{Editorial scope.} Counted unit: the article's lemma lem:selector (source0.tex lines 141--154). The article's complete original proof of it is delivered with it (source0.tex lines 156--275, one \texttt{proof} environment of 120 lines). No mathematics is written, repaired or re-derived: the statement and the proof are the article's own lines, in the article's own notation, and no retained line is substituted or re-flowed.  No further source-local context is needed: every cross-reference of every retained line resolves inside the two delivered spans.  Nothing was omitted from any retained span, so the package carries no omission record.  The retained lines cite 1 works, whose entries are transcribed verbatim from the article's own bibliography source in the e-print and delivered with short editorial labels recorded in the item's reference mapping.  No retained line contains a figure environment, an image-inclusion command or drawing code, so no image is delivered and the package declares no asset dependency.  Editorial formatting only, in addition: an AMS \texttt{amsart} preamble that restates the article's own declarations and macros used by the retained lines, namely the environments \texttt{lemma} on separate counters, together with the standard packages those lines need.  The retained environments print in this document as Lemma 1 in order of appearance, whereas the article prints the same items with the numbering of its own full document.  The complete native source of the article as distributed in the arXiv e-print for this exact version is bundled unchanged as \texttt{sources/207538/source0.tex}.
\begin{lemma}[Contact-selector lemma]\label{lem:selector}
Let $F:\Hh\to\D$ be univalent, and let $E\subset\R$ be a finite union of
pairwise disjoint open intervals, possibly unbounded. Assume that
\begin{enumerate}[label=\textup{(\roman*)},leftmargin=2.4em]
\item for almost every $x>0$, exactly one of $x,-x$ belongs to $E$;
\item $F$ extends holomorphically across every interval comprising $E$;
\item $|F|=1$ on $E$.
\end{enumerate}
Then
\begin{equation}\label{eq:selector-main}
  \int_0^\infty |F'(iy)|\,dy
  \le \frac\pi2\int_E|F'(x)|\,dx.
\end{equation}
\end{lemma}

\begin{proof}
We may assume
\begin{equation}\label{eq:Afinite}
  A:=\int_E|F'(x)|\,dx<\infty.
\end{equation}

\smallskip
\noindent\emph{Factorization and boundary positivity.}
The function $F$ has at most one zero. If $F(a)=0$, where $a=u+iv$ and
$v>0$, put
\[
  B=B_a,\qquad B_a(z)=\frac{z-a}{z-\overline a};
\]
otherwise put $B\equiv1$. Schwarz--Pick gives $|F|\le|B_a|$ in the first
case. Hence in either case
\begin{equation}\label{eq:FBQ}
  F=BQ,
  \qquad Q\ne0,
  \qquad |B|,|Q|\le1.
\end{equation}
Write $Q=e^{ih}$. Then $\operatorname{Im}h\ge0$, and either $h$ is real
constant or $h:\Hh\to\Hh$ is a Pick function.

On every interval contained in $E$, the function $Q$ extends holomorphically,
is nonvanishing, and is unimodular. A local logarithm therefore extends $h$
holomorphically through that interval, with real boundary values. The Herglotz representation \cite[Chapter~I]{Duren} has the form
\begin{equation}\label{eq:Herglotz}
  h(z)=c+\alpha z+
  \int_{\R}\left(\frac1{t-z}-\frac{t}{1+t^2}\right)d\mu(t),
\end{equation}
where $c\in\R$, $\alpha\ge0$, and $\mu$ is positive with
$\int(1+t^2)^{-1}\,d\mu(t)<\infty$. By the standard Stieltjes inversion
formula, holomorphic continuation with real boundary values through an
interval implies that $\mu$ has no mass there. Consequently, for $x\in E$,
\begin{equation}\label{eq:hprime-boundary}
  h'(x)=\alpha+\int_{\R}\frac{d\mu(t)}{(t-x)^2}\ge0.
\end{equation}

If the zero $a=u+iv$ is present, set
\[
  p(x)=\frac{2v}{(x-u)^2+v^2};
\]
otherwise set $p=0$. Since $B'(x)/B(x)=ip(x)$ on $\R$, we obtain the exact
identity
\begin{equation}\label{eq:no-cancellation}
  |F'(x)|=p(x)+h'(x),
  \qquad x\in E.
\end{equation}
The selector has infinite Lebesgue measure. Thus \eqref{eq:Afinite},
\eqref{eq:hprime-boundary}, and \eqref{eq:no-cancellation} force
\begin{equation}\label{eq:alpha-zero}
  \alpha=0.
\end{equation}

\smallskip
\noindent\emph{The zero-free factor.}
Schwarz--Pick for $h$ and \eqref{eq:Herglotz} give
\[
  |Q'(iy)|
  \le\frac{\operatorname{Im}h(iy)}{y}
  =\int_{\R}\frac{d\mu(t)}{t^2+y^2}.
\]
Hence, by Tonelli,
\begin{equation}\label{eq:Qvertical}
  \int_0^\infty|Q'(iy)|\,dy
  \le\frac\pi2\int_{\R}\frac{d\mu(t)}{|t|}.
\end{equation}
On the other hand, the selector property gives, for every $t\ne0$,
\begin{equation}\label{eq:kernel}
  \int_E\frac{dx}{(x-t)^2}\ge\frac1{|t|}.
\end{equation}
Indeed, for $t>0$ each pair $x,-x$ contributes one of
$(x-t)^{-2},(x+t)^{-2}$, and the smaller one integrates over $x>0$ to $1/t$;
the case $t<0$ is symmetric, while for $t=0$ the inner integral is infinite.
Therefore
\begin{equation}\label{eq:Qestimate}
  \int_Eh'(x)\,dx
  =\int_{\R}\left(\int_E\frac{dx}{(x-t)^2}\right)d\mu(t)
  \ge\int_{\R}\frac{d\mu(t)}{|t|},
\end{equation}
and
\begin{equation}\label{eq:Qfinal}
  \int_0^\infty|Q'(iy)|\,dy
  \le\frac\pi2\int_Eh'(x)\,dx.
\end{equation}

\smallskip
\noindent\emph{The possible zero.}
If $a=u+iv$ is present, then
\[
  |B_a'(iy)|=\frac{2v}{u^2+(y+v)^2}.
\]
For $u\ne0$, writing $r=v/|u|$, we have
\begin{equation}\label{eq:Bvertical}
  \int_0^\infty|B_a'(iy)|\,dy=2r\arctan(1/r),
\end{equation}
whereas the selector property yields
\begin{equation}\label{eq:Bboundary}
  \int_Ep(x)\,dx
  \ge\int_0^\infty\min\{p(x),p(-x)\}\,dx
  =2\arctan r.
\end{equation}
The elementary inequality
\[
  2r\arctan(1/r)\le\pi\arctan r
\]
follows on writing $r=\tan s$ and using
$\cos s\ge2(\pi/2-s)/\pi$ and $\sin s\le s$, both for
$0\leq s \leq \pi/2$. Hence
\begin{equation}\label{eq:Bfinal}
  \int_0^\infty|B_a'(iy)|\,dy
  \le\frac\pi2\int_Ep(x)\,dx.
\end{equation}
For $u=0$, the two sides before multiplication by $\pi/2$ are respectively
$2$ and $\pi$, so \eqref{eq:Bfinal} still holds.

Finally, $F=BQ$ and $|B|,|Q|\le1$ imply $|F'|\le|B'|+|Q'|$ on $i\R_+$.
Combining \eqref{eq:Qfinal}, \eqref{eq:Bfinal}, and
\eqref{eq:no-cancellation} proves \eqref{eq:selector-main}.
\end{proof}

\begin{thebibliography}{99}
\bibitem[Duren]{Duren} P. L. Duren, \emph{Theory of $H^p$ Spaces}, Pure and Applied Mathematics, vol. 38, Academic Press, New York, 1970.
\end{thebibliography}
\end{document}
